\documentclass[12pt,a4paper]{article}
\usepackage{amsmath}
\usepackage{amssymb}
\usepackage{graphicx}
\usepackage{float}
\usepackage{booktabs}
\usepackage{hyperref}
\usepackage{geometry}
\geometry{margin=1in}

\title{ENSO Forecasting: MIFS Feature Selection and XGBoost Implementation}
\author{}
\date{\today}

\begin{document}

\maketitle

\section*{IMPLEMENTATION}

\subsection*{1. Mutual Information Feature Selection (MIFS)}

MIFS is a greedy feature selection algorithm that balances relevance (mutual information with target) and redundancy (mutual information among selected features).

\subsubsection*{1.1 Algorithm Overview}

The MIFS score for feature $X_{ij}$ is defined as:
\begin{equation}
S(X_{ij}) = I(X_{ij}; Y) - \frac{1}{|S|} \sum_{X_{kl} \in S} I(X_{ij}; X_{kl})
\end{equation}

where:
\begin{itemize}
    \item $I(X_{ij}; Y)$ is the mutual information (relevance) between feature and target
    \item $I(X_{ij}; X_{kl})$ is the mutual information (redundancy) between feature pair
    \item $S$ is the current set of selected features
\end{itemize}

\subsubsection*{1.2 Implementation Steps}

\paragraph{Step 1: Compute Relevance}
\begin{itemize}
    \item For all $\sim 20,000$ features in windowed dataset
    \item Use Kraskov k-NN estimator (sklearn's \texttt{mutual\_info\_regression})
    \item k-neighbors parameter: $k=5$ for MI computation
    \item Training set only: 457 samples (window end $\leq$ 2018-12)
\end{itemize}

\paragraph{Step 2: Stratified Pre-filtering}
Unlike flat top-N selection, stratified prefiltering uses per-variable-type quotas:
\begin{itemize}
    \item Group features by variable type (e.g., \texttt{sst\_anom\_z}, \texttt{sohtc300\_anom\_z})
    \item Per-type quota = $\min(\text{cap\_max}, \max(\text{cap\_min}, \text{fraction} \times n_{\text{type}}))$
    \item Parameters: fraction = 0.4, cap\_min = 15, cap\_max = 30
    \item Result: Candidate pool of $\leq 300$ features with all variable types represented
    \item Rationale: Prevents SST (strongest individual predictor) from monopolizing candidate pool; ensures mechanistically important but weak predictors (OHC, wind stress) are included
\end{itemize}

\paragraph{Step 3: Greedy MIFS Loop}
\begin{itemize}
    \item Initialize: $S = \emptyset$
    \item For each step $t = 1, \ldots, \min(50, |\text{candidates}|)$:
    \begin{enumerate}
        \item For each remaining candidate $X_i$:
        \item If $|S| = 0$: score $= I(X_i; Y)$ (pure relevance)
        \item Else: score $= I(X_i; Y) - \frac{1}{|S|} \sum_{X_j \in S} I(X_i; X_j)$ (redundancy penalty)
        \item Cache pairwise MI to avoid recomputation
    \end{enumerate}
    \item Select: $X^* = \arg\max_i \text{score}(X_i)$
    \item Update: $S \leftarrow S \cup \{X^*\}$
    \item Output: Ranked list of selected features in order
\end{itemize}

\paragraph{Step 4: K-sweep with Validation}
\begin{itemize}
    \item For each $K = 1, 2, \ldots, 50$ in steps of 2:
    \item Take top-$K$ features from greedy ranking
    \item Train XGBoost (baseline: $n\_\text{est}=200$, $\text{max\_depth}=3$, $\text{lr}=0.1$) on train set (457 samples)
    \item Evaluate on validation set (48 samples, window end 2019-01 to 2022-12)
    \item Compute RMSE on validation predictions
    \item Select: $K^* = \arg\min_K \text{RMSE}_{\text{val}}(K)$
\end{itemize}

\begin{figure}[H]
\centering
\includegraphics[width=0.75\textwidth]{Outputs/pipeline_output_tuned/mifs_ksweep_plot.png}
\caption{K-sweep validation: RMSE vs. number of features K. The optimal K minimizes validation error.}
\label{fig:ksweep}
\end{figure}

\paragraph{Step 5: Output}
Save:
\begin{itemize}
    \item Full relevance ranking (all $\sim 20,000$ features)
    \item MIFS-selected features in ranking order
    \item K-sweep results (RMSE for each K)
    \item Best K and final feature list
\end{itemize}

---

\subsection*{2. XGBoost Regression}

XGBoost (eXtreme Gradient Boosting) is an ensemble tree model that sequentially adds trees to minimize loss.

\subsubsection*{2.1 Regularization-Focused Grid Search}

To prevent overfitting observed in earlier runs, a **regularized grid** was designed:

\begin{table}[H]
\centering
\begin{tabular}{lll}
\toprule
\textbf{Parameter} & \textbf{Baseline} & \textbf{Regularized} \\
\midrule
max\_depth & 6 & \textbf{2, 3} \\
n\_estimators & 500 & 100, 200, 300 \\
learning\_rate & 0.1 & 0.01, 0.05, 0.1 \\
subsample & — & 0.6, 0.7, 0.8 \\
colsample\_bytree & — & 0.6, 0.7, 0.8 \\
min\_child\_weight & — & 5, 10, 20 \\
reg\_alpha (L1) & — & 0, 0.5, 1.0 \\
reg\_lambda (L2) & — & 1, 5, 10 \\
\bottomrule
\end{tabular}
\caption{Regularized XGBoost grid (4,374 combinations)}
\end{table}

\subsubsection*{2.2 Hyperparameter Tuning Process}

\paragraph{Stage 1: Grid Search on Validation Set}
\begin{itemize}
    \item Enumerate all parameter combinations ($3 \times 3 \times 3 \times 3 \times 3 \times 3 \times 3 \times 3 = 6561$ when all expanded)
    \item Actually: $2 \times 3 \times 3 \times 3 \times 3 \times 3 \times 3 \times 3 = 4374$ combos (max\_depth $\in \{2, 3\}$)
    \item For each combo: Train on train set only (457 samples), evaluate on validation set (48 samples)
    \item Metric: Validation RMSE
    \item Select best hyperparameters: $\theta^* = \arg\min_\theta \text{RMSE}_{\text{val}}(\theta)$
\end{itemize}

\paragraph{Stage 2: Final Evaluation on Test Set}
\begin{itemize}
    \item Re-train model with $\theta^*$ on train set only (457 samples)
    \item Evaluate on held-out test set (35 samples, window end 2023-01 to 2023-12+)
    \item Metrics: RMSE, ACC (anomaly correlation coefficient)
    \item Note: Validation set (48 samples) is held out and not used in final training
\end{itemize}

\subsubsection*{2.3 Anomaly Correlation Coefficient (ACC)}

For ENSO forecasting, ACC is more operationally relevant than raw RMSE:
\begin{equation}
\text{ACC} = r(y_{\text{true}}, y_{\text{pred}}) = \frac{\text{Cov}(y_{\text{true}}, y_{\text{pred}})}{\sigma_{\text{true}} \cdot \sigma_{\text{pred}}}
\end{equation}

\begin{itemize}
    \item Pearson correlation between predicted and actual Niño 3.4 index
    \item Range: $[-1, 1]$; higher is better
    \item Measures whether model gets the sign and phase correct
\end{itemize}

\paragraph{Why ACC Matters for ENSO:}
\begin{itemize}
    \item Forecast skill depends on predicting El Niño/La Niña \textbf{phase} correctly
    \item A prediction 0.05°C off but correct sign $>$ a prediction 0.01°C off but wrong sign
\end{itemize}

---

\newpage

\section*{RESULT AND ANALYSIS}

\subsection*{1. MIFS Feature Selection Results}

\subsubsection*{1.1 Stratified vs. Non-Stratified Prefiltering}

\begin{table}[H]
\centering
\begin{tabular}{llll}
\toprule
\textbf{Model} & \textbf{Prefilter Type} & \textbf{Best K} & \textbf{Top Features} \\
\midrule
Baseline & Flat top-300 & 22 & sossheig, t2m, sohtc300, so20chgt \\
Stratified & Quota-based & 10 & sossheig, t2m, sst, msl \\
Tuned & Tuned (flat) & 18 & sossheig, sohtc300, t2m, so20chgt \\
Regularized & Tuned (flat) & 18 & (same as Tuned) \\
Stratified+Reg & Quota-based & 10 & sossheig, t2m, sst, msl, avg\_inss \\
\bottomrule
\end{tabular}
\caption{MIFS feature counts and top predictors across models}
\end{table}

\paragraph{Key Observations:}
\begin{itemize}
    \item \textbf{Flat prefiltering} selects 18-30 features; typically dominated by strongest individual predictors (SST, OHC)
    \item \textbf{Stratified prefiltering} selects fewer (10) but guarantees per-type representation
    \item \textbf{Top-ranked feature across all runs}: \texttt{sossheig\_anom\_z\_(0.0,192.0)\_lag0} (sea surface height at Niño 3.4 box, current month)
    \item \textbf{Consistent predictors}: Ocean heat content (sohtc300), thermocline (so20chgt), atmospheric (t2m, msl)
\end{itemize}

\subsubsection*{1.2 Redundancy Penalty Effect}

Why does MIFS select so few features from $\sim 20,000$?

\begin{itemize}
    \item \textbf{Spatial correlation}: All grid cells in same ocean basin are highly correlated
    \item Once one SST cell is selected, all nearby cells get large redundancy penalty
    \item MIFS greedy loop effectively clusters spatially-correlated features and picks one representative
    \item With only 457 training samples, adding more features causes validation RMSE to increase (curse of dimensionality)
    \item K-sweep finds sweet spot: typically 10-22 features minimize validation error
\end{itemize}

---

\subsection*{2. XGBoost Performance Across Models}

\subsubsection*{2.1 Baseline vs. Tuned vs. Regularized}

\begin{table}[H]
\centering
\begin{tabular}{llllll}
\toprule
\textbf{Model} & \textbf{Features} & \textbf{Train RMSE} & \textbf{Val RMSE} & \textbf{Test RMSE} & \textbf{Test ACC} \\
\midrule
Baseline & 22 & not logged & 0.6821 & 0.8554 & 0.5608 \\
Stratified & 10 & not logged & 0.7089 & 0.8949 & 0.4972 \\
Tuned & 18 & \textbf{0.0005} (!) & 0.6054 & 0.8906 & 0.8078 \\
Regularized & 18 & 0.3934 $\checkmark$ & 0.6000 & 0.9077 & \textbf{0.8163} $\checkmark$ \\
Stratified+Reg & 10 & 0.3581 $\checkmark$ & 0.6534 & 0.9466 & 0.5872 \\
\bottomrule
\end{tabular}
\caption{Model performance comparison (Train/Val/Test RMSE and ACC)}
\end{table}

\subsubsection*{2.2 Key Results}

\paragraph{Baseline (Flat MIFS + Standard XGB):}
\begin{itemize}
    \item 22 selected features
    \item Test RMSE: 0.8554°C (predictions off by ~0.86°C on average)
    \item Test ACC: 0.5608 (moderate skill at predicting sign)
    \item No train/val data logged; unclear generalization
\end{itemize}

\paragraph{Tuned (Tuned MIFS + Expanded XGB Grid):}
\begin{itemize}
    \item 18 selected features
    \item \textbf{Severe overfitting}:
    \begin{itemize}
        \item Train RMSE = 0.0005 (essentially perfect on training data)
        \item Train ACC = 1.00 (correlates perfectly with training targets)
        \item Test RMSE = 0.8906 (generalizes poorly)
        \item Test ACC = 0.8078 (good on test despite overfitting)
    \end{itemize}
    \item Train-test gap: +0.8901 RMSE (model memorized training set)
    \item \textbf{Not recommended for operational use}
\end{itemize}

\begin{figure}[H]
\centering
\includegraphics[width=0.8\textwidth]{Outputs/pipeline_output_tuned/test_predictions_plot.png}
\caption{Tuned model test predictions: actual vs. predicted Niño 3.4 index (2023-2025).}
\label{fig:tuned_predictions}
\end{figure}

\begin{figure}[H]
\centering
\includegraphics[width=0.75\textwidth]{Outputs/pipeline_output_tuned/feature_importances_plot.png}
\caption{Tuned model feature importances: shows which features XGBoost relied on most.}
\label{fig:tuned_importance}
\end{figure}

\paragraph{Regularized (Tuned MIFS + Regularized XGB Grid):}
\begin{itemize}
    \item 18 selected features (same as Tuned, using reused MIFS from Tuned)
    \item \textbf{Controlled overfitting}:
    \begin{itemize}
        \item Train RMSE = 0.3934 (learning, not memorizing)
        \item Train ACC = 0.9139 (strong but not perfect)
        \item Test RMSE = 0.9077 (only +0.517 worse than train)
        \item Test ACC = \textbf{0.8163} (best among all models)
    \end{itemize}
    \item Train-test gap: +0.5143 RMSE (reasonable, indicates good generalization)
    \item Regularization parameters:
    \begin{itemize}
        \item max\_depth = 2 (shallow trees only)
        \item min\_child\_weight = 20 (high leaf threshold)
        \item reg\_lambda = 5 (strong L2 penalty)
        \item subsample = 0.6, colsample\_bytree = 0.6 (aggressive subsampling)
    \end{itemize}
    \item \textbf{Recommended for operational deployment}
\end{itemize}

\begin{figure}[H]
\centering
\includegraphics[width=0.8\textwidth]{Outputs/pipeline_output_regularized/overfitting_comparison.png}
\caption{Regularized model overfitting analysis: train vs. validation performance across epochs.}
\label{fig:overfitting}
\end{figure}

\begin{figure}[H]
\centering
\includegraphics[width=0.8\textwidth]{Outputs/pipeline_output_regularized/test_predictions_plot.png}
\caption{Regularized model test predictions: actual vs. predicted Niño 3.4 index (2023-2025). Higher correlation than tuned model.}
\label{fig:reg_predictions}
\end{figure}

\begin{figure}[H]
\centering
\includegraphics[width=0.75\textwidth]{Outputs/pipeline_output_regularized/feature_importances_plot.png}
\caption{Regularized model feature importances: top predictors from stratified MIFS selection.}
\label{fig:reg_importance}
\end{figure}

\paragraph{Stratified+Regularized (Quota-based MIFS + Regularized XGB Grid):}
\begin{itemize}
    \item \textbf{10 selected features} (smallest model)
    \item Combines stratified quota-based MIFS with anti-overfitting regularization
    \item \textbf{Performance metrics}:
    \begin{itemize}
        \item Train RMSE = 0.3581 (controlled learning)
        \item Train ACC = 0.9297 (strong)
        \item Val RMSE = 0.6534
        \item Val ACC = 0.4038 (lower than expected)
        \item Test RMSE = 0.9466 (highest among all models)
        \item Test ACC = 0.5872 (moderate)
    \end{itemize}
    \item Train-test gap: +0.5885 RMSE (good generalization)
    \item Best hyperparameters:
    \begin{itemize}
        \item max\_depth = 2 (very shallow)
        \item min\_child\_weight = 20 (high leaf threshold)
        \item reg\_lambda = 5 (strong L2 penalty)
        \item subsample = 0.6, colsample\_bytree = 0.8
    \end{itemize}
    \item \textbf{Trade-off}: Most concise model (only 10 features) with good generalization, but lower test ACC than Regularized (18 features)
\end{itemize}

\subsubsection*{2.3 Model Comparison Summary}

\begin{table}[H]
\centering
\begin{tabular}{ll}
\toprule
\textbf{Metric} & \textbf{Winner} \\
\midrule
Best Test RMSE & Baseline (0.8554) \\
Best Test ACC & Regularized (0.8163) \\
Fewest Features & Stratified+Reg \& Stratified (10) \\
Best for Deployment & \textbf{Regularized} \\
\bottomrule
\end{tabular}
\caption{Recommended model: Regularized (best balance of RMSE, ACC, and robustness)}
\end{table}

---

\subsection*{3. Analysis}

\subsubsection*{3.1 Why Regularized Outperforms Tuned}

\textbf{The Tuned model problem}:
\begin{equation}
\text{Train RMSE} = 0.0005 \Rightarrow \text{Model has memorized every training sample}
\end{equation}

When training data is only 457 samples and model has $\sim$100 features $\times$ 6 levels (boosting depth), the model has enough capacity to fit every point perfectly. This is useless for forecasting new months.

\textbf{The Regularized model solution}:
\begin{itemize}
    \item \textbf{Shallow trees (max\_depth=2)}: Cannot fit complex, high-dimensional patterns
    \item \textbf{High min\_child\_weight (20)}: Prevents leaf nodes for single samples
    \item \textbf{L2 penalty (reg\_lambda=5)}: Shrinks leaf weight magnitudes
    \item \textbf{Row/column subsampling}: Each tree sees only 60\% of rows/features
    \item Train RMSE stays $\approx 0.39$ (real learning, not memorization)
    \item Test performance degrades gracefully to 0.91 RMSE (plausible generalization)
\end{itemize}

\subsubsection*{3.2 Anomaly Correlation vs. RMSE}

For operational ENSO forecasting, **ACC often matters more than raw RMSE**:

\begin{itemize}
    \item \textbf{Baseline}: Lower RMSE (0.8554) but moderate ACC (0.5608)
    \begin{itemize}
        \item Can predict magnitude but not always the phase (El Niño vs. La Niña)
    \end{itemize}
    \item \textbf{Regularized}: Slightly higher RMSE (0.9077) but highest ACC (0.8163)
    \begin{itemize}
        \item Gets the sign/phase right in 81.6\% of predictions
        \item More useful for seasonal forecasting
    \end{itemize}
\end{itemize}

\paragraph{Example:}
\begin{itemize}
    \item True Niño 3.4 = $+1.5°$C (El Niño)
    \item Baseline predicts: $+1.48°$C (RMSE = 0.02, but if internally correlated poorly, ACC low)
    \item Regularized predicts: $+1.42°$C (RMSE = 0.08, but correlation with true is strong, ACC high)
    \item For forecasting, Regularized is more reliable
\end{itemize}

\subsubsection*{3.3 Stratified vs. Flat MIFS Trade-off}

\begin{itemize}
    \item Flat (Baseline, Tuned, Regularized): Selects 18-22 features, mostly SST + OHC
    \item Stratified (Stratified, Stratified+Reg): Selects 10 features, guarantees all variable types
    \begin{itemize}
        \item More concise model (fewer parameters)
        \item More physically interpretable (all known ENSO predictors present)
        \item Includes weaker-but-important predictors: avg\_inss (north-south wind stress)
    \end{itemize}
    \item \textbf{Performance comparison}:
    \begin{itemize}
        \item Stratified (no regularization): Test RMSE=0.8949, ACC=0.4972
        \item Stratified+Reg (with regularization): Test RMSE=0.9466, ACC=0.5872
        \item Regularized (flat MIFS, 18 features): Test RMSE=0.9077, ACC=0.8163
    \end{itemize}
    \item \textbf{Key finding}: Stratified MIFS with only 10 features achieves moderate performance (ACC=0.59) but falls short of flat MIFS + regularization (ACC=0.82)
    \item \textbf{Interpretation}: The highly correlated SST grid cells selected by flat MIFS carry more predictive signal than the quota-enforced diversity
\end{itemize}

\subsubsection*{3.4 Feature Count vs. Performance}

\begin{itemize}
    \item \textbf{10 features (Stratified+Regularized)}: ACC=0.59 — too few, missing key patterns
    \item \textbf{18 features (Regularized)}: ACC=0.82 — sweet spot for 457 training samples
    \item \textbf{22 features (Baseline)}: ACC=0.56 — without regularization, starts overfitting
\end{itemize}


\end{document}
